% !TEX program = xelatex
\documentclass[aspectratio=169,10pt]{beamer}

\input{preamble2}
% part2-defs.tex -- part 2 macros and styles (Proteek Rosul)
%   loaded by part2.tex and by ../combined.tex

% ======================================================================
%  3D SUPPORT
% ======================================================================
\usepackage{pgfplots}
\pgfplotsset{compat=1.18}

% overlay-aware visibility for tikz / pgfplots elements:
%   \addplot3[..., visible on=<2->] ...      \draw[visible on=<3->] ...
\tikzset{
  invisible/.style  = {opacity=0, text opacity=0},
  visible on/.style = {alt={#1{}{invisible}}},
  alt/.code args    = {<#1>#2#3}{\alt<#1>{\pgfkeysalso{#2}}{\pgfkeysalso{#3}}},
}

% height ramps: deep at the goal, pale on the rim.  Blue = Euclidean,
% orange = Manhattan -- the same pairing the 2D distance slide uses.
\pgfplotsset{
  colormap={hblue}{color=(rbluedk) color=(rbluemd) color=(rbluelt)},
  colormap={horange}{color=(cbOrange!70!black) color=(cbOrange) color=(cbOrangeLt)},
  % flat grid tiles, value = 0 free / 1 expanded / 2 start / 3 goal
  colormap={htiles}{color=(gfree) color=(gclose) color=(gstart) color=(ggoal)},
  land3d/.style   = {axis lines=none, clip=false},
  % land view={az}{el}{unit in cm}: pgfplots' own view={az}{el} projection,
  % but with one fixed length per data unit, so shapes keep their true
  % proportions (a cone stays round) and the size is predictable
  land view/.code n args = {3}{%
    \pgfmathsetmacro\lvxx{#3*cos(#1)}%
    \pgfmathsetmacro\lvxy{-#3*sin(#1)*sin(#2)}%
    \pgfmathsetmacro\lvyx{#3*sin(#1)}%
    \pgfmathsetmacro\lvyy{#3*cos(#1)*sin(#2)}%
    \pgfmathsetmacro\lvzy{#3*cos(#2)}%
    \edef\lvkeys{\noexpand\pgfkeysalso{%
      x={(\lvxx cm,\lvxy cm)}, y={(\lvyx cm,\lvyy cm)}, z={(0cm,\lvzy cm)}}}%
    \lvkeys},
  % lifted grid: flat-shaded facets with white gutters, like the 2D tiles
  tilesurf/.style = {surf, shader=faceted, faceted color=white,
                     line width=0.45pt, z buffer=sort, point meta=z},
}

\tikzset{
  sball/.style  = {circle, fill=gstart, draw=white, line width=0.7pt,
                   minimum size=3.4mm, inner sep=0pt},
  gball/.style  = {sball, fill=ggoal},
  vball/.style  = {sball, fill=hlyellow, draw=cbOrange, line width=1pt},
  tag/.style    = {font=\scriptsize, inner sep=1.5pt},
  tagbox/.style = {tag, fill=white, fill opacity=0.85, text opacity=1,
                   rounded corners=1pt},
  move3d/.style = {-{Stealth[length=2.2mm]}, line width=1.7pt,
                   line cap=round, shorten >=0.6pt},
}

% ---------- per-figure geometry ----------------------------------------
% 3D-a: 12 x 12 lattice, G at (7,3), height = kA * Euclidean distance
\def\kA{0.45}\def\gAx{7}\def\gAy{3}
\newcommand{\liftA}[2]{(#1,#2,{\kA*veclen((#1)-\gAx,(#2)-\gAy)})}
% the flat map on the floor (a loop body with # cannot sit inside a frame)
\newcommand{\floorA}{%
  \pgfplotsinvokeforeach{0,...,12}{%
    \draw[gridgray, line width=0.35pt] (##1,0,0) -- (##1,12,0);
    \draw[gridgray, line width=0.35pt] (0,##1,0) -- (12,##1,0);}}

% 3D-b: G at the origin, v at (6,0), height = kB * Euclidean distance
\def\kB{0.45}
\newcommand{\liftB}[2]{(#1,#2,{\kB*veclen(#1,#2)})}

% 3D-c: 18 x 12 cells, S and G cell centres, optimal cost C* = 5
%       Dijkstra expands the disk  |vS| <= C*;
%       A* the ellipse             |vS| + |vG| <= C* + slack
\def\kC{0.3}
\def\cSx{5.5}\def\cSy{6.5}\def\cGx{10.5}\def\cGy{6.5}
\def\cCost{5}\def\cSlack{0.9}
\newcommand{\legendC}{%
  \begin{tikzpicture}[x=0.34cm,y=0.34cm, baseline=0pt]
    \renewcommand{\gpad}{0.09}%
    \gcell{gclose}{0}{0}
      \node[anchor=west, font=\scriptsize] at (1.4,0.5) {expanded};
    \gcell{gfree}{7}{0}
      \node[anchor=west, font=\scriptsize] at (8.4,0.5) {never touched};
    \gcell{cbGreen}{16.5}{0}
      \node[anchor=west, font=\scriptsize] at (17.9,0.5) {optimal route};
  \end{tikzpicture}}

% 3D-e: cone / pyramid, rim at distance 4, sample point P = (2,1)
\def\kE{0.8}\def\pEx{2}\def\pEy{1}


% colour-coded highlight, shown from overlay 2 on: #1 = 2 orange, 3 green,
% 4 blue, 5 red, anything else yellow.  preamble2.tex calls it \hlon; it
% is renamed so it can sit beside part 1's \hlon{<overlay>}{...} in
% combined.tex.
\newcommand{\hlonc}[2]{%
  \alt<2->{%
    \setlength{\fboxsep}{1.4pt}%
    \ifnum#1=2 \colorbox{cbOrange!30}{#2}%
    \else\ifnum#1=3 \colorbox{cbGreen!30}{#2}%
    \else\ifnum#1=4 \colorbox{rblue!30}{#2}%
    \else\ifnum#1=5 \colorbox{cbRed!30}{#2}%
    \else \colorbox{hlyellow}{#2}%
    \fi\fi\fi\fi
  }{#2}%
}
          % part 2 macros (also loaded by ../combined.tex)
\input{grids/stats.tex}

\begin{document}

% ---------- shared macro: a small labelled grid + a distance readout ---
\newcommand{\distgrid}[3]{%   #1 = draw commands (tikz), #2 = caption, #3 = formula
  \begin{minipage}[t]{4.55cm}
    \centering
    \begin{tikzpicture}[x=0.34cm,y=0.34cm]
      \foreach \c in {0,...,8}{\foreach \r in {0,...,5}{\gcell{gfree}{\c}{\r}}}
      #1
    \end{tikzpicture}\par
    \vspace{2pt}
    {\footnotesize #2}\par
    {\scriptsize\mut{#3}}
  \end{minipage}}

 
% ---------- 11a. what f, g, h are ----------------------------------------
\begin{frame}[label=11a]{Defining parameters}
  \begin{columns}[T,onlytextwidth]
  \column{0.5\textwidth}
    \onslide<2->{%
    \begin{block}{$g(v)$ --- cost so far}
      \small
      The exact cost of the \textbf{best known path} from the start $s$ to $v$.
    \end{block}}
    \vspace{0.5em}
    \onslide<3->{%
    \begin{exampleblock}{$h(v)$ --- cost to go}
      \small
      An \textbf{estimate} of the remaining cost from $v$ to the goal $t$.
    \end{exampleblock}}
    \onslide<4->{%
    \begin{alertblock}{$f(v) = g(v) + h(v)$ --- total estimated cost}
      \small
      \textbf{Best guess} for the cost of the \emph{full} $s \to t$ path if it routes through $v$.
    \end{alertblock}}

  \column{0.4\textwidth}
    \vspace{0.6em}

    \onslide<5->{%
    \begin{tikzpicture}[x=0.42cm,y=0.42cm]
      \foreach \c in {0,...,8}{\foreach \r in {0,...,4}{\gcell{gfree}{\c}{\r}}}
      \gletter{gstart}{0}{2}{S}
      \gletter{ggoal}{8}{2}{G}
      \draw[cbGreen, line width=1.6pt, line cap=round]
        (0.7,2.5) -- (3.3,2.5);
      \gletter{gclose}{3}{2}{v}
      \draw[cbOrange, dashed, line width=1.6pt, line cap=round]
      (3.8,2.5) -- (8.2,2.5);
      {\node[font=\tiny, text=cbGreen] at (2,2) {$g(v)$};}
      {\node[font=\tiny, text=cbOrange] at (6,2) {$h(v)$};}
    \end{tikzpicture}\par
    \vspace{4pt}
    {\footnotesize\textcolor{cbGreen}{$g(v)=3$} --- three steps walked}\par
    {\footnotesize\textcolor{cbOrange}{$h(v)=5$} --- five steps remaining}\par
    {\footnotesize $f(v) = g(v) + h(v) = 3 + 5 = 8$}
    }
  \end{columns}
\end{frame}
 
 
% ---------- 12a. why the heuristic helps --------------------------------
\begin{frame}{Why the heuristic helps}
  \begin{columns}[T]
  \column{0.5\textwidth}
    \centering
    \onslide<1->{%
    {\footnotesize\textbf{No heuristic} --- Dijkstra ($h \equiv 0$)}\\[2pt]
    \gridpicsmall{0.34cm}{grids/mini-dijkstra}\par
    {\scriptsize\textcolor{cbRed}{\statdijkstraexp} cells expanded --- \emph{blind} to where the goal is}}

  \column{0.5\textwidth}
    \centering
    \onslide<2->{%
    {\footnotesize\textbf{With a good heuristic} --- A*}\\[2pt]
    \gridpicsmall{0.34cm}{grids/mini-astar}\par
    {\scriptsize\textcolor{cbGreen}{\statastarexp} cells expanded ---
     \emph{biased} toward the goal}}
  \end{columns}

  \vspace{0.6em}
  \onslide<3->{\takeaway{$h$ does not change \emph{what} A* is allowed to
    explore --- it changes the \emph{order}.}}
\end{frame}


% ---------- 3D-a. the heuristic as a landscape -------------------------
\begin{frame}{The heuristic as a landscape}
  \begin{columns}[T]
  \column{0.58\textwidth}
    \centering
    \begin{tikzpicture}
      \begin{axis}[land3d, land view={-32}{34}{0.38}]
        % the flat map, left on the floor
        \floorA
        % the same map, every node lifted to height h(v)
        \addplot3[tilesurf, colormap name=hblue, samples=13, samples y=13,
                  domain=0:12, y domain=0:12]
          {\kA*veclen(x-\gAx, y-\gAy)};
        % height dimension at the near corner
        \draw[neutralg, line width=0.6pt, {Bar[width=1.6mm]}-{Bar[width=1.6mm]}]
          (12,0,0) -- \liftA{12}{0}
          node[tagbox, text=neutralg, midway, left=2pt] {$h(v)$};
        % the route, rolling downhill
        \draw[edgpath, line cap=round, line join=round, visible on=<2->]
          \liftA{2}{10} -- \liftA{3}{9} -- \liftA{4}{8} -- \liftA{5}{7}
          -- \liftA{6}{6} -- \liftA{7}{5} -- \liftA{7}{4} -- \liftA{7}{3};
        % one step the wrong way
        \draw[move3d,cbRed, visible on=<3->] \liftA{2}{10} -- \liftA{1}{11};
        \node[tagbox, text=cbRed, above=3pt, visible on=<3->]
          at \liftA{1}{11} {uphill: $\Delta f=+2$};
        \node[tagbox, text=cbGreen, left=5pt, visible on=<3->]
          at \liftA{4}{8} {downhill: $\Delta f=0$};
        \node[sball] at \liftA{2}{10} {};
        \node[tag, text=gstart, below left=3pt] at \liftA{2}{10} {\textbf{S}};
        \node[gball] at \liftA{7}{3} {};
        \node[tag, text=ggoal, below=6pt] at \liftA{7}{3} {\textbf{G} \ $h=0$};
      \end{axis}
    \end{tikzpicture}

  \column{0.42\textwidth}
    \begin{block}{Lift the map by $h$}
      \small
      \begin{itemize}\setlength{\itemsep}{5pt}
        \item<1-> raise every node $v$ to height $h(v)$
        \item<1-> $h(G)=0$
        \item<2-> each step costs $+1$ in $g$
        \item<3-> a \pos{downhill} step is incentivized
        \item<3-> an \bad{uphill} step is penalized
      \end{itemize}
    \end{block}
  \end{columns}
\end{frame}


% ---------- 12b. heuristic choice, expanded ---------------------------
\begin{frame}[label=12b]{Choosing a heuristic for A*}
  \begin{columns}[T]
  \column{0.52\textwidth}
    \begin{block}{What $h(v)$ has to be}
      \small
      \begin{itemize}\setlength{\itemsep}{7pt}
        \item<2-> Cheap to compute
        \item<3-> \textbf{Never an overestimate}
        \item<4-> Ideally close to the true remaining cost
      \end{itemize}
    \end{block}

  \column{0.48\textwidth}
    \onslide<5->{%
    \begin{exampleblock}{Common choices}
      \small
      \begin{itemize}\setlength{\itemsep}{6pt}
        \item<5-> Manhattan distance
        \item<5-> Euclidean distance
      \end{itemize}
    \end{exampleblock}}
  \end{columns}
\end{frame}

% ---------- 12c. Manhattan vs Euclidean, diagrams ----------------------
\begin{frame}{Manhattan vs Euclidean distance}
  \centering
  \begin{columns}[T]
  \column{0.5\textwidth}
    \centering
    \onslide<1->{%
    \distgrid{%
      \gletter{gstart}{1}{1}{A}%
      \gletter{ggoal}{6}{4}{B}%
      \onslide<1->{\draw[cbOrange, line width=1.6pt, line cap=round]
        (1.8,1.5) -- (6.5,1.5);}%
      \onslide<1->{\draw[cbOrange, line width=1.6pt, line cap=round]
        (6.5,1.5) -- (6.5,4.2);}%
      \onslide<1->{\node[font=\tiny, text=cbOrange] at (4.0,1.0) {5};}%
      \onslide<1->{\node[font=\tiny, text=cbOrange] at (6.9,3.0) {3};}%
    }{\textbf{Manhattan} \\ $|x_A - x_B| + |y_A - y_B|$}
     {only horizontal / vertical steps --- $5+3=8$}}
  \column{0.5\textwidth}
    \centering
    \onslide<2->{%
    \distgrid{%
      \gletter{gstart}{1}{1}{A}%
      \gletter{ggoal}{6}{4}{B}%
      \draw[rblue, line width=1.6pt, line cap=round] (1.7,1.7) -- (6.3,4.3);
      \node[font=\tiny, text=rblue] at (4.7,2.7) {5.8};
    }{\textbf{Euclidean} \\ $\sqrt{(x_A-x_B)^2+(y_A-y_B)^2}$}
     {any-angle movement --- $\sqrt{5^2+3^2}\approx 5.8$}}
  \end{columns}
\end{frame}

% ---------- 3D-d. Euclidean vs Manhattan in 3D ---------------------------
\begin{frame}{Euclidean vs Manhattan --- in three dimensions}
  \begin{columns}[T]
  \column{0.55\textwidth}
    \centering
    \begin{tikzpicture}[x={(1.22cm,0cm)}, y={(0.66cm,0.48cm)},
                        z={(0cm,1.22cm)}, line join=round, line cap=round]
      % faint cell lattice on the floor
      \foreach \i in {0,...,4}{\draw[gridgray, line width=0.4pt] (\i,0,0) -- (\i,3,0);}
      \foreach \j in {0,...,3}{\draw[gridgray, line width=0.4pt] (0,\j,0) -- (4,\j,0);}
      % the hidden edge
      \draw[gridgray, dashed, line width=0.6pt] (0,3,0) -- (0,3,2);
      % Pythagoras triangle (floor diagonal, rise, space diagonal)
      \fill[rbluelt, opacity=0.8, visible on=<2->] (0,0,0) -- (4,3,0) -- (4,3,2) -- cycle;
      \draw[rbluemd, dashed, line width=1pt, visible on=<2->] (0,0,0) -- (4,3,0)
        node[pos=0.6, below right=-1pt, tag, text=rbluemd] {$5$};
      % the box
      \fill[rbluelt, opacity=0.35] (0,0,2) -- (4,0,2) -- (4,3,2) -- (0,3,2) -- cycle;
      \draw[rbluedk!70, line width=0.8pt]
        (0,0,0) -- (4,0,0) -- (4,0,2) -- (0,0,2) -- cycle
        (4,0,0) -- (4,3,0) -- (4,3,2) -- (4,0,2)
        (0,0,2) -- (0,3,2) -- (4,3,2);
      % Manhattan: one axis at a time, along the edges
      \draw[cbOrange, line width=2.2pt, visible on=<1->]
        (0,0,0) -- (4,0,0) -- (4,3,0) -- (4,3,2);
      \node[tag, text=cbOrange, below=3pt, visible on=<1->] at (2,0,0) {$\Delta x = 4$};
      \node[tag, text=cbOrange, below right=0pt, visible on=<1->] at (4,1.5,0) {$\Delta y = 3$};
      \node[tagbox, text=cbOrange, left=3pt, visible on=<1->] at (4,3,1) {$\Delta z = 2$};
      % Euclidean: straight through the box
      \draw[rblue, line width=2.2pt, visible on=<2->] (0,0,0) -- (4,3,2);
      \node[tagbox, text=rblue, above left=0pt, visible on=<2->] at (2,1.5,1)
        {$\sqrt{29}\approx 5.39$};
      % endpoints
      \node[sball] at (0,0,0) {};
      \node[tag, text=gstart, left=6pt] at (0,0,0) {\textbf{A}};
      \node[gball] at (4,3,2) {};
      \node[tag, text=ggoal, above=6pt] at (4,3,2) {\textbf{B}};
    \end{tikzpicture}

  \column{0.45\textwidth}
    \onslide<1->{%
    \begin{block}{\textcolor{cbOrange}{Manhattan} --- along the edges}
      \small
      $|\Delta x|+|\Delta y|+|\Delta z| = 4+3+2 = \mathbf{9}$\\[2pt]
      {\footnotesize one axis at a time, like a 6-directional grid}
    \end{block}}
    \onslide<2->{%
    \begin{block}{\textcolor{rblue}{Euclidean} --- straight through}
      \small
      $\sqrt{\Delta x^2+\Delta y^2+\Delta z^2} = \sqrt{29} \approx \mathbf{5.39}$\\[2pt]
      {\footnotesize any-angle movement}
    \end{block}}
    \vspace{0.3em}
    \vspace{0.3em}
    \onslide<3->{{\footnotesize Manhattan is \textcolor{cbRed}{never smaller}, equal only
      when $A$ and $B$ differ along a single axis.}}
  \end{columns}
\end{frame}


% ---------- 12d. pseudocode for A* --------------------------------------
\begin{frame}[t]{A* --- pseudocode}
  \begin{columns}[T]

    \column{0.55\textwidth}
    \begin{clrs}[\footnotesize]
      \aline[1-]{\proc{A-Star}$(G, s, t, h)$}
      \aline[1-]{\idt1 \kw{for each} $v \in V$: \
        $g[v] \gets \infty$,\; $\pi[v] \gets \textsc{nil}$}
      \aline[1-]{\idt1 $g[s] \gets 0$}
      \aline[1-]{\idt1 \hlonc{2}{$Q \gets$ priority queue keyed by
        \textcolor{cbGreen}{$g$}$+$\textcolor{cbOrange}{$h$}, holding
        $\{s\}$}}
      \aline[1-]{\idt1 $\mathit{S} \gets \varnothing$}
      \aline[1-]{\idt1 \kw{while} $Q \neq \varnothing$}
      \aline[1-]{\idt2 $u \gets \textsc{Extract-Min}(Q)$}
      \aline[1-]{\idt2 \kw{if} $u = t$:}
      \aline[1-]{\idt3 \kw{return} \proc{Reconstruct}$(\pi, t)$}
      \aline[1-]{\idt2 $\mathit{S} \gets \mathit{S} \cup \{u\}$}
      \aline[1-]{\idt2 \kw{for each} $v$ adjacent to $u$}
      \aline[1-]{\idt3 \kw{if} $v \in \mathit{S}$: \ \kw{continue}}
      \aline[1-]{\idt3 \kw{if} $g[u] + w(u,v) < g[v]$}
      \aline[1-]{\idt4 $g[v] \gets g[u] + w(u,v)$,\; $\pi[v] \gets u$}
      \aline[1-]{\idt4 \hlonc{2}{\textsc{Insert-or-Decrease-Key}
        $\big(Q, v, g[v]+h(v)\big)$}}
    \end{clrs}

    \column{0.45\textwidth}
    \onslide<2->{%
    \begin{block}{Only one difference from Dijkstra}
      \small The priority key changes from $g$ to $(g+h)$. Everywhere
      else, the algorithm is untouched.
    \end{block}}
  \end{columns}
\end{frame}

% ---------- 12e. time complexity, annotated pseudocode ---------------------
\begin{frame}[t, label=12e]{Time complexity analysis}
  \begin{columns}[T]

    \column{0.55\textwidth}
    \begin{clrs}[\footnotesize]
      \aline[1-]{\proc{A-Star}$(G, s, t, h)$}
      \aline[1-]{\idt1 \hlonc{2}{\kw{for each} $v \in V$: \
        $g[v] \gets \infty$,\; $\pi[v] \gets \textsc{nil}$}}
      \aline[1-]{\idt1 $g[s] \gets 0$}
      \aline[1-]{\idt1 $Q \gets$ priority queue keyed by $g+h$, holding
        $\{s\}$}
      \aline[1-]{\idt1 $\mathit{S} \gets \varnothing$}
      \aline[1-]{\idt1 \hlonc{3}{\kw{while} $Q \neq \varnothing$}}
      \aline[1-]{\idt2 \hlonc{3}{$u \gets \textsc{Extract-Min}(Q)$}}
      \aline[1-]{\idt2 \kw{if} $u = t$:}
      \aline[1-]{\idt3 \kw{return} \proc{Reconstruct}$(\pi, t)$}
      \aline[1-]{\idt2 $\mathit{S} \gets \mathit{S} \cup \{u\}$}
      \aline[1-]{\idt2 \hlonc{4}{\kw{for each} $v$ adjacent to $u$}}
      \aline[1-]{\idt3 \kw{if} $v \in \mathit{S}$: \ \kw{continue}}
      \aline[1-]{\idt3 \kw{if} $g[u] + w(u,v) < g[v]$}
      \aline[1-]{\idt4 $g[v] \gets g[u] + w(u,v)$,\;
        $\pi[v] \gets u$}
      \aline[1-]{\idt4 \hlonc{5}{\textsc{Insert-or-Decrease-Key}
        $\big(Q, v, g[v]+h(v)\big)$}}
    \end{clrs}

    \column{0.45\textwidth}
    \onslide<2->{%
    \begin{exampleblock}{\textcolor{cbOrange}{Init}}
      \scriptsize $O(|V|)$ for touching every vertex once
    \end{exampleblock}
    \vspace{0.2em}
    \begin{exampleblock}{\textcolor{cbGreen}{Main loop}}
      \scriptsize $O(|V|\log|V|)$ on a binary heap
    \end{exampleblock}
    \vspace{0.2em}
    \begin{exampleblock}{\textcolor{rblue}{Edge scan}}
      \scriptsize $O(|E|)$ for considering every edge
    \end{exampleblock}
    \vspace{0.2em}
    \begin{alertblock}{\textcolor{cbRed}{Heap update}}
      \scriptsize $O(\log|V|)$ each, up to $O(|E|\log|V|)$ total
    \end{alertblock}}
  \end{columns}
\end{frame}

% ---------- 12e-2. time complexity, the sum -------------------------------
\begin{frame}[label=12e2]{Time complexity --- putting it together}
  % \small
  \begin{itemize}\setlength{\itemsep}{9pt}
    \item<1-> \textbf{\textcolor{cbOrange}{Initialization}}: $O(|V|)$
    \item<2-> \textbf{\textcolor{cbGreen}{Extractions}}: $O(|V|\log|V|)$
    \item<3-> \textbf{\textcolor{rblue}{Edge scans}}: $O(|E|)$
    \item<4-> \textbf{\textcolor{cbRed}{Heap insert / decrease-key}}: $O(|E|\log|V|)$
  \end{itemize}

  \vspace{0.6em}
  \onslide<5->{%
  \[
    O(|V|) + O(|V|\log|V|) + O(|E|) + O(|E|\log|V|)
  \]}
  \onslide<6->{%
  \begin{center}
    \Large\(\Rightarrow\quad
    \boxed{O\big((|V|+|E|)\log |V|\big)}\)
  \end{center}}
\end{frame}

% ---------- 12e-3. space complexity, annotated pseudocode ------------------
\begin{frame}[t, label=12e3]{Space complexity analysis}
  \begin{columns}[T]

    \column{0.55\textwidth}
    \begin{clrs}[\footnotesize]
      \aline[1-]{\proc{A-Star}$(G, s, t, h)$}
      \aline[1-]{\idt1 \hlonc{2}{\kw{for each} $v \in V$: \
        $g[v] \gets \infty$,\; $\pi[v] \gets \textsc{nil}$}}
      \aline[1-]{\idt1 $g[s] \gets 0$}
      \aline[1-]{\idt1 \hlonc{3}{$Q \gets$ priority queue keyed by
        $g+h$, holding $\{s\}$}}
      \aline[1-]{\idt1 \hlonc{3}{$\mathit{S} \gets \varnothing$}}
      \aline[1-]{\idt1 \kw{while} $Q \neq \varnothing$}
      \aline[1-]{\idt2 $u \gets \textsc{Extract-Min}(Q)$}
      \aline[1-]{\idt2 \kw{if} $u = t$:}
      \aline[1-]{\idt3 \hlonc{4}{\kw{return} \proc{Reconstruct}$(\pi, t)$}}
      \aline[1-]{\idt2 \hlonc{3}{$\mathit{S} \gets \mathit{S} \cup \{u\}$}}
      \aline[1-]{\idt2 \hlonc{3}{\kw{for each} $v$ adjacent to $u$}}
      \aline[1-]{\idt3 \kw{if} $v \in \mathit{S}$: \ \kw{continue}}
      \aline[1-]{\idt3 \kw{if} $g[u] + w(u,v) < g[v]$}
      \aline[1-]{\idt4 \hlonc{3}{$g[v] \gets g[u] + w(u,v)$,\;
        $\pi[v] \gets u$}}
      \aline[1-]{\idt4 \hlonc{3}{\textsc{Insert-or-Decrease-Key}
        $\big(Q, v, g[v]+h(v)\big)$}}
    \end{clrs}

    \column{0.45\textwidth}
    \onslide<2->{%
    \begin{exampleblock}{\textcolor{cbOrange}{$g$ and $\pi$ arrays}}
      \scriptsize One entry per vertex: $O(|V|)$
    \end{exampleblock}
    \vspace{0.4em}
    \begin{exampleblock}{\textcolor{cbGreen}{Frontier $Q$ and closed set $S$}}
      \scriptsize Every vertex is generated at most once: together $O(|V|)$.
    \end{exampleblock}
    \vspace{0.4em}
    \begin{alertblock}{\textcolor{rblue}{Reconstructed path}}
      \scriptsize $O(|V|)$ as typically $O(d) \ll |V|$
    \end{alertblock}}
  \end{columns}
\end{frame}

% ---------- 12e-4. space complexity, the sum --------------------------------
\begin{frame}[label=12e4]{Space complexity --- putting it together}
  \small
  \begin{itemize}\setlength{\itemsep}{9pt}
    \item<1-> \textbf{\textcolor{cbOrange}{$g[\cdot]$ and
          $\pi[\cdot]$}}: $O(|V|)$
    \item<2-> \textbf{\textcolor{cbGreen}{Priority queue $Q$ and
          closed set $\mathit{S}$}}: $O(|V|)$
    \item<3-> \textbf{\textcolor{rblue}{Path reconstruction}}: $O(|V|)$ worst case
  \end{itemize}

  \vspace{0.6em}
  \onslide<4->{%
  \[
    O(|V|) + O(|V|) + O(|V|)
  \]}
  \onslide<5->{%
  \begin{center}
    \Large\(\Rightarrow\quad \boxed{O(|V|)}\)
  \end{center}}
  \onslide<6->{%
  \begin{center}
    \mut{No dependence on $|E|$.} }
  \end{center}
\end{frame}

% ---------- 12f. complexity comparison table -----------------------------
\begin{frame}[label=12f]{Complexity, side by side}
  \centering
  {\footnotesize
  \renewcommand{\arraystretch}{1.18}
  \begin{tabular}{@{}l@{\hspace{1.4em}}c@{\hspace{1.4em}}c@{\hspace{1.4em}}c@{}}
    \toprule
     & \textcolor{rbluemd}{\shortstack{\textbf{Dijkstra}}} & \textcolor{cbOrange}{\shortstack{\textbf{Greedy}}}
     & \textcolor{cbGreen}{\shortstack{\textbf{A*}}}\\
    \midrule
    {\shortstack[l]{Time (worst case)}}
      & {$O((|V|{+}|E|)\log|V|)$}
      & {$O((|V|{+}|E|)\log|V|)$}
      & {$O((|V|{+}|E|)\log|V|)$}\\[5pt]
    {Space}
      & {$O(|V|)$} & {$O(|V|)$} & {$O(|V|)$}\\[3pt]
    {\shortstack[l]{Nodes expanded}}
      & {most / all} & {few, unpredictable}
      & {few, provably safe}\\[5pt]
    {Optimal?}
      & {\textcolor{cbGreen}{$\checkmark$}} & {\textcolor{cbRed}{$\times$}}
      & {\textcolor{cbGreen}{$\checkmark$}}\\
    \bottomrule
  \end{tabular}\par}

  \vspace{0.8em}
  \onslide<2->{\begin{center}\footnotesize\mut{Same asymptotic worst
    case as Dijkstra --- A* wins in \emph{practice}, not in the
    textbook bound.}\end{center}}
\end{frame}

\begin{frame}[label=12g0]{Correctness conditions}
  \begin{center}
  \begin{tikzpicture}
    \node<1->[draw=rblue, fill=rbluelt, rounded corners=3pt, line width=1pt,
        text width=4.8cm, align=center, inner sep=12pt] at (0,0)
      {{\small\mut{1}}\\[3pt]{\bfseries\large\HL{Admissibility}}};
    \node<2->[draw=cbOrange, fill=cbOrangeLt, rounded corners=3pt,
        line width=1pt, text width=4.8cm, align=center, inner sep=12pt]
        at (6.2,0)
      {{\small\mut{2}}\\[3pt]{\bfseries\large\con{Consistency}}};
  \end{tikzpicture}
  \end{center}
\end{frame}

% ---------- 12g. admissibility ------------------------------------------
\begin{frame}[label=12g]{Correctness --- admissibility}
  \begin{columns}[T]
  \column{0.55\textwidth}
    \begin{block}{Definition}
      \small A heuristic $h$ is \textbf{admissible} if, for every node
      $v$,
      \[ h(v) \;\le\; h^{*}(v) \]
      where $h^{*}(v)$ is the \emph{true} cheapest cost from $v$ to the
      goal.
    \end{block}
    \vspace{0.6em}
    \onslide<2->{%
    \begin{exampleblock}{Why it guarantees optimality}
      \small When A* pops $t$, every other node in $Q$ has
      $f = g+h \ge$ that $f$-value. Since $h$ never overestimates,
      nothing left in the queue can beat $t$'s true cost.
    \end{exampleblock}}

  \column{0.45\textwidth}
    \onslide<2->{%
    \centering
    \begin{tikzpicture}[x=0.85cm,y=0.85cm]
      \node[vstart] (v) at (0,0) {$v$};
      \node[vgoal]  (t) at (3,0.3) {$t$};
      \draw[edgpath, {Stealth[length=2mm]}-, -{Stealth[length=2mm]}]
        (v) to[bend left=18] node[wlab, above] {$h^{*}(v)=6$} (t);
      \draw[edgorange, dashed, -{Stealth[length=2mm]}]
        (v) to[bend right=18] node[wlab, below] {$h(v)=4$} (t);
    \end{tikzpicture}\par
    {\footnotesize\textcolor{cbGreen}{$h(v)=4 \le h^{*}(v)=6$} ---
     admissible}\par
    \vspace{4pt}
    {\footnotesize\textcolor{cbRed}{if instead $h(v)=8$} --- not
     admissible}}
  \end{columns}
\end{frame}

% ---------- 12h. consistency --------------------------------------------
\begin{frame}[label=12h]{Correctness --- consistency}
  \begin{columns}[T]
  \column{0.55\textwidth}
    \begin{block}{Definition}
      \small A heuristic $h$ is \textbf{consistent} if, for
      every edge $(u,v)$ with weight $w(u,v)$,
      \[ h(u) \;\le\; w(u,v) + h(v) \]
      with $h(\text{goal})=0$.
    \end{block}
    \vspace{0.6em}
    \onslide<2->{%
    \begin{exampleblock}{Why it matters}
      \small The \textbf{triangle inequality} applied to $h$: \\ $f=g+h$ never
      decreases along a path.
    \end{exampleblock}}

  \column{0.45\textwidth}
    \onslide<2->{%
    \centering
    \begin{tikzpicture}[x=0.85cm,y=0.85cm]
      \node[vstart] (u) at (0,1.2)  {$u$};
      \node[vclosed]  (v) at (3,0)    {$v$};
      \node[vgoal]  (t) at (1.5,-1.8) {$G$};

      \draw[edgpath, -{Stealth[length=2mm]}]
        (u) -- node[wlab, above, yshift=8pt] {$w(u,v)$} (v);
      \draw[edgpath, -{Stealth[length=2mm]}]
        (v) -- node[wlab, right, xshift=8pt] {$h(v)$} (t);
      \draw[edgorange, dashed, -{Stealth[length=2mm]}]
        (u) -- node[wlab, left, xshift=-8pt] {$h(u)$} (t);
    \end{tikzpicture}\par
    {\footnotesize \mut{$h(u) \le w(u,v) + h(v)$ --- the triangle inequality in effect.}}}

  \end{columns}
\end{frame}

% ---------- 12j. a bad heuristic, worked examples (expanded, animated) ---
\begin{frame}{A bad heuristic (I) --- overestimating}
  \centering
  {\footnotesize\textbf{Overestimating $h$} --- looks like Greedy}\\[4pt]
  \begin{tikzpicture}[x=0.42cm,y=0.42cm]
    \onslide<1->{%
    \foreach \c in {0,...,8}{\foreach \r in {0,...,5}{\gcell{gfree}{\c}{\r}}}
    \foreach \r in {0,1,2,3}{\gcell{gwall}{4}{\r}}
    \gletter{gstart}{0}{2}{S}
    \gletter{ggoal}{8}{2}{G}}
    \onslide<2->{\draw[cbOrange, line width=1.6pt, line cap=round]
      (0.5,2.8) -- (0.5,4.5);}
    \onslide<3->{\draw[cbOrange, line width=1.6pt, line cap=round]
      (0.5,4.5) -- (4.5,4.5);}
    \onslide<4->{\draw[cbRed, line width=1.6pt, line cap=round]
      (4.5,4.5) -- (4.5,2.5);
      \node[font=\tiny, text=cbRed] at (5.6,4.9) {rushes at the wall};}
    \onslide<5->{\draw[cbRed, line width=1.6pt, line cap=round]
      (4.5,2.5) -- (8.2,2.5);}
  \end{tikzpicture}\par
  \vspace{2pt}
  \begin{itemize}
    \item<5-> \footnotesize Inadmissible
    \item<5-> \footnotesize Behaves like Greedy
    \item<5-> \footnotesize \textcolor{cbRed}{Can return a longer-than-optimal path}
  \end{itemize}
\end{frame}

% ---------- 12j-2. uninformative heuristic --------------------------------
\begin{frame}{A bad heuristic (II) --- uninformative}
  \centering
  {\footnotesize\textbf{Uninformative $h$} --- looks like Dijkstra}\\[4pt]
  \begin{tikzpicture}[x=0.42cm,y=0.42cm]
    \onslide<1->{%
    \foreach \c in {0,...,8}{\foreach \r in {0,...,5}{\gcell{gfree}{\c}{\r}}}
    \gletter{gstart}{0}{2}{S}
    \gletter{ggoal}{8}{2}{G}}
    \onslide<2->{\foreach \c in {0,1,2}{\foreach \r in {0,2,4}{\gcell{gclose}{\c}{\r}}}}
    \onslide<3->{\foreach \c in {3,4,5}{\foreach \r in {0,2,4}{\gcell{gclose}{\c}{\r}}}}
    \onslide<4->{\foreach \c in {6,7,8}{\foreach \r in {0,2,4}{\gcell{gclose}{\c}{\r}}}}
  \end{tikzpicture}\par
  \vspace{2pt}
  \begin{itemize}
    \item<4-> \footnotesize $h(v)=1$ for all non-goal $v$
    \item<4-> \footnotesize \textbf{Admissible}, but slow --- same as Dijkstra
  \end{itemize}
\end{frame}

% ---------- 12j-3. admissible but locally misleading ----------------------
\begin{frame}{A bad heuristic (III) --- admissible but misleading}
  \centering
  {\footnotesize\textbf{Admissible $h$, dead-end pocket} --- correct,
   yet still wasteful}\\[4pt]
   \vspace{0.6em}
  \begin{tikzpicture}[x=0.42cm,y=0.42cm]
    \onslide<1->{%
    \foreach \c in {0,...,8}{\foreach \r in {0,...,5}{\gcell{gfree}{\c}{\r}}}
    \foreach \r in {3,4,5}{\gcell{gwall}{3}{\r}}
    \foreach \r in {3,4,5}{\gcell{gwall}{5}{\r}}
    \gcell{gwall}{4}{5}
    \gletter{gstart}{0}{2}{S}
    \gletter{ggoal}{8}{5}{G}}
    \onslide<2->{\draw[cbOrange, line width=1.6pt, line cap=round]
      (0.8,2.5) -- (4.5,2.5);}
    \onslide<3->{\draw[cbRed, line width=1.6pt, line cap=round]
      (4.5,2.5) -- (4.5,4.7);
      \node[font=\tiny, text=cbRed] at (4.5,6.5) {dead end};}
    \onslide<4->{\draw[cbRed, line width=1.3pt, line cap=round]
      (4.5, 4.7) -- (4.7,4.7) -- (4.7,2.5);}
    \onslide<5->{\draw[cbGreen, line width=1.6pt, line cap=round]
      (4.7,2.5) -- (8.5,2.5) -- (8.5,5.2);}
  \end{tikzpicture}\par
  \vspace{2pt}
  \begin{itemize}
    \item<5-> \footnotesize Still \textbf{admissible} --- just blind to the obstacle
    \item<5-> \footnotesize \bad{Finds the optimal path, but wastes a detour} on a dead end
  \end{itemize}
\end{frame}

% ---------- 12j-4. three failure modes, side by side -----------------------
\begin{frame}[label=12j4]{Three ways a heuristic can go wrong}
  \centering
  \begin{tikzpicture}[
      algbox/.style = {draw, line width=0.8pt, align=center, font=\small,
                       minimum width=4.7cm, minimum height=1.85cm, inner sep=5pt},
      flow/.style   = {-{Stealth[length=2.4mm]}, line width=1.0pt, draw=neutralg},
      flowlab/.style= {font=\footnotesize, text=neutralg, inner sep=2pt}]

    \node<2->[algbox, draw=cbRed, fill=cbRedLt] (large) at (-4.2,2)
      {\textcolor{cbRed}{\textbf{Too large}}\\[3pt]
       \footnotesize inadmissible --- overestimates \\ the \footnotesize remaining cost \\[5pt]
       \textcolor{cbRed}{$\times$}~correctness\qquad
       \textcolor{cbGreen}{$\checkmark$}~decisive};

    \node<3->[algbox, draw=neutralg, fill=neutralg!12] (flat) at (2.2,2)
      {\textcolor{neutralg}{\textbf{Too flat}}\\[3pt]
       \footnotesize uninformative --- admissible \\ \footnotesize but barely varies \\[5pt]
       \textcolor{cbGreen}{$\checkmark$}~correct\qquad
       \textcolor{cbRed}{$\times$}~degrades to Dijkstra};

    \node<4->[algbox, draw=cbOrange, fill=cbOrangeLt] (blind) at (-1.2,-1)
      {\textcolor{cbOrange}{\textbf{Blind to structure}}\\[3pt]
       \footnotesize (admissible, locally misleading) --- \\
       \footnotesize ignores walls and dead ends\\[5pt]
       \textcolor{cbGreen}{$\checkmark$}~correct\qquad
       \textcolor{cbRed}{$\times$}~wastes expansions};

  \end{tikzpicture}
\end{frame}

\end{document}